% GENERATED from content/data/enrichment.json by mxbuild.enrichment; do not edit.
\mxunit{enrichment-telescoping-solutions}{Sommes télescopiques}
\mxsubtitle{Première--Terminale · suites \textperiodcentered{} Solutions détaillées \textperiodcentered{} 10}

\begin{exercise}
Calculer $\sum [\frac{1}{k}-\frac{1}{k+1}]$ de $k=3$ à 10.
\end{exercise}
\begin{solution}
Les termes intermédiaires s’annulent : $\frac{1}{3}-\frac{1}{11}=\frac{8}{33}$.
\end{solution}

\begin{exercise}
Décomposer $\frac{1}{k(k+2)}$ en une différence de deux fractions.
\end{exercise}
\begin{solution}
$\frac{1}{k(k+2)}=(\frac{1}{2})[\frac{1}{k}-\frac{1}{k+2}]$.
\end{solution}

\begin{exercise}
Calculer $\sum (2k+1)$ de $k=0$ à n en utilisant $(k+1)^{2}-k^{2}$.
\end{exercise}
\begin{solution}
$2k+1=(k+1)^{2}-k^{2}$. La somme vaut $(n+1)^{2}-0^{2}=(n+1)^{2}$.
\end{solution}

\begin{exercise}
Si $S_{n}=\frac{n}{n+1}$, trouver le plus petit entier $n\geq 1$ tel que $1-S_{n}<0{,}01$.
\end{exercise}
\begin{solution}
$\frac{1}{n+1}<\frac{1}{100}$ équivaut à $n+1>100$, donc $n\geq 100$. Le minimum est 100.
\end{solution}

\begin{problem}[Chercher, justifier, conclure]
Soit $T_{n}=\sum  \frac{1}{k(k+2)}$ pour k de 1 à n. Établir une expression exacte en fonction de n, puis calculer sa limite. Trouver le plus petit n tel que l’écart à la limite soit strictement inférieur à $0{,}01$.
\end{problem}
\begin{solution}
Après décomposition, tous les termes sauf les deux premiers et les deux derniers s’annulent : $T_{n}=\frac{3}{4}-\frac{1}{2(n+1)}-\frac{1}{2(n+2)}$. La limite est $\frac{3}{4}$. L’écart décroît avec n. Pour $n=98$, il vaut $\frac{1}{198}+\frac{1}{200}>0{,}01$ ; pour $n=99$, il vaut $\frac{1}{200}+\frac{1}{202}<0{,}01$. Le minimum est donc 99.
\end{solution}

